\documentclass[10pt]{article}
\usepackage[margin=1in]{geometry}
\usepackage{fontspec}
\setmainfont{texgyretermes}[Extension=.otf, UprightFont=*-regular, BoldFont=*-bold, ItalicFont=*-italic, BoldItalicFont=*-bolditalic]
\usepackage{booktabs}
\usepackage{caption}
\title{Supplement: LLM Chart Judges Reward Labels, Audiences Reward Topics}
\date{}
\begin{document}
\maketitle
All tables are generated by \texttt{src/11\_supplement.py} from the CSVs in \texttt{analysis/}. Pairs and prompts are identical to the main paper unless stated.

\section*{S1. Data cleaning}
Filters in order (\texttt{src/02\_clean.py}). Deviations from our original plan:
(i) a new filter for blank images (grayscale SD $<$ 2), which are video thumbnails labeled from their titles;
(ii) the perceptual-hash duplicate threshold was tightened from Hamming 4 to 2, because visual QA showed that distance-4 matches were different charts with shared layouts;
(iii) a new rule drops groups of 3 or more near-duplicates whose titles are unrelated (median token-set similarity $<$ 60), which removes link previews and site banners;
(iv) transparent images are composited on white.
No rule uses votes.

\begin{table}[h]\centering\small
\caption{Cleaning funnel.}
\begin{tabular}{lrr}
\toprule
Step & Rows remaining & Removed \\
\midrule
raw rows & 52836 & 0 \\
unreadable image, missing score or date & 52836 & 0 \\
blank image (grayscale SD < 2) & 52688 & 148 \\
chart type empty / none / Other & 46239 & 6449 \\
short side < 300 px & 44433 & 1806 \\
score < 5 & 35928 & 8505 \\
near-duplicate (pHash <= 2) or placeholder group & 34556 & 1372 \\
\bottomrule
\end{tabular}
\end{table}

\section*{S2. Other prompts, title-only control, close pairs}
P2 is a six-dimension VisJudge-style rubric (1--5 each, averaged). P3 is a pairwise prompt run in both orders: credit 1 if both orders pick the winner, 0 if both pick the loser, 0.5 if they disagree. The title-only control asks gpt-5 which of two titles got more upvotes, in both orders. These rows come from the earlier, smaller runs: P3 on the first 250 or 500 clear pairs, P2 on the first 150, and close pairs (vote ratio 1.5--2.5) on 200. $\Delta$ is the judge minus the image-blind baseline on the same pairs (95\% bootstrap CI).

\begin{table}[h]\centering\small
\caption{Rubric (P2), pairwise (P3), title-only, and close-pair results.}
\begin{tabular}{llrrrrrrr}
\toprule
Set & Row & $n$ & Acc. & $\Delta$ & lo & hi & Order flips & Picks first \\
\midrule
clear & Gemini 2.5 Flash P3 & 500 & 0.564 & -0.042 & -0.089 & +0.007 & 0.53 & 0.77 \\
clear & gpt-5 title only & 500 & 0.582 & -0.024 & -0.074 & +0.024 & 0.50 & 0.26 \\
clear & gpt-5 P3 & 250 & 0.636 & +0.052 & -0.028 & +0.126 & 0.18 & 0.54 \\
clear & Claude Haiku 4.5 P3 & 250 & 0.564 & -0.020 & -0.102 & +0.064 & 0.18 & 0.42 \\
clear & gpt-5 P2 & 150 & 0.600 & +0.080 & -0.023 & +0.183 & -- & -- \\
close & Gemini 2.5 Flash P1 & 200 & 0.512 & +0.022 & -0.068 & +0.113 & -- & -- \\
clear & feature baseline (OOF) & 500 & 0.606 & -- & -- & -- & -- & -- \\
close & feature baseline (OOF) & 200 & 0.490 & -- & -- & -- & -- & -- \\
\bottomrule
\end{tabular}
\end{table}

\begin{table}[h]\centering\small
\caption{Prompt contrasts. The rubric and pairwise formats do not beat holistic P1, and images do not beat titles for gpt-5.}
\begin{tabular}{lrrrr}
\toprule
Paired contrast (same pairs) & $n$ & Diff. & lo & hi \\
\midrule
gpt-5 P2 minus gpt-5 P1 & 150 & -0.043 & -0.103 & +0.013 \\
gpt-5 P3 minus gpt-5 P1 & 250 & +0.000 & -0.048 & +0.046 \\
Claude Haiku 4.5 P3 minus Claude Haiku 4.5 P1 & 250 & -0.020 & -0.080 & +0.032 \\
Gemini 2.5 Flash P3 minus Gemini 2.5 Flash P1 & 500 & -0.023 & -0.054 & +0.006 \\
gpt-5 P3 minus gpt-5 title only & 250 & +0.008 & -0.062 & +0.074 \\
\bottomrule
\end{tabular}
\end{table}

\begin{table}[h]\centering\small
\caption{Robustness check: P1 equivalence on the 1,500 pairs no judge saw during development.}
\begin{tabular}{lrrrrrl}
\toprule
Comparison (P1, 1,500 fresh pairs) & $n$ & Diff. & 90\% lo & 90\% hi & TOST $p$ & Verdict \\
\midrule
gpt-5 P1 minus baseline & 1500 & -0.0277 & -0.0553 & +0.0000 & 0.233 & inconclusive \\
Claude Haiku 4.5 P1 minus baseline & 1500 & -0.0317 & -0.0587 & -0.0047 & 0.305 & inconclusive \\
Gemini 2.5 Flash P1 minus baseline & 1500 & -0.0140 & -0.0403 & +0.0127 & 0.055 & inconclusive \\
\bottomrule
\end{tabular}
\end{table}

\section*{S3. Calibration on VisJudge-Bench}
150 single-view items sampled with seed 20261008. P2 rubric mean vs.\ expert \texttt{overall\_score}. Published values (VisJudge-Bench v3, 648-item test split over all types): GPT-5 $r=0.428$, MAE 0.553. The test split cannot be reconstructed from the released file, so the comparison is approximate.

\begin{table}[h]\centering\small
\caption{P2 calibration against expert ratings.}
\begin{tabular}{lrrrrr}
\toprule
Judge & $n$ & Pearson $r$ & lo & hi & MAE \\
\midrule
Gemini 2.5 Flash & 150 & 0.472 & 0.319 & 0.599 & 0.871 \\
gpt-5 & 150 & 0.414 & 0.262 & 0.537 & 0.622 \\
Claude Haiku 4.5 & 150 & 0.467 & 0.326 & 0.590 & 0.692 \\
\bottomrule
\end{tabular}
\end{table}

\begin{table}[h]\centering\small
\caption{Same prompt (P1), same metric (Spearman), different targets. Sonnet 5.5 was not run on calibration items because of the budget cap.}
\begin{tabular}{lrrrrr}
\toprule
Judge & P1 $\rho$ vs.\ experts & Audience images & P1 $\rho$ vs.\ audience pct & lo & hi \\
\midrule
gpt-5 & 0.442 & 4000 & 0.084 & 0.051 & 0.112 \\
Claude Haiku 4.5 & 0.489 & 3999 & 0.115 & 0.086 & 0.145 \\
Gemini 2.5 Flash & 0.453 & 4000 & 0.133 & 0.103 & 0.164 \\
Claude Sonnet 5.5 & -- & 2000 & 0.119 & 0.077 & 0.164 \\
\bottomrule
\end{tabular}
\end{table}

\section*{S4. Memorization}
Training cutoffs: gpt-5 2024-09-30, Claude Haiku 4.5 Feb 2025 (reliable), Gemini 2.5 Flash Jan 2025. The dataset ends Jan 2025, so only gpt-5 has post-cutoff pairs. The recognition probe asked each judge whether it had seen each of the 50 most-voted charts and, if so, for the title. Because 72\% of these charts print their title, recalls that match the chart's own OCR text are not counted as memory.

\begin{table}[h]\centering\small
\caption{Accuracy by training-cutoff group (first 500 clear pairs).}
\begin{tabular}{lllrrrr}
\toprule
Judge & Prompt & Cutoff group & $n$ & Acc. & lo & hi \\
\midrule
gpt-5 & P1 & post & 37 & 0.581 & 0.446 & 0.703 \\
gpt-5 & P1 & pre & 463 & 0.579 & 0.537 & 0.619 \\
gpt-5 & P3 & post & 17 & 0.706 & 0.500 & 0.882 \\
gpt-5 & P3 & pre & 233 & 0.631 & 0.569 & 0.687 \\
Claude Haiku 4.5 & P1 & pre & 500 & 0.555 & 0.520 & 0.590 \\
Claude Haiku 4.5 & P3 & pre & 250 & 0.564 & 0.512 & 0.620 \\
Gemini 2.5 Flash & P1 & pre & 500 & 0.587 & 0.555 & 0.618 \\
Gemini 2.5 Flash & P3 & pre & 500 & 0.564 & 0.535 & 0.593 \\
\bottomrule
\end{tabular}
\end{table}

\begin{table}[h]\centering\small
\caption{Recognition probe.}
\begin{tabular}{lrrrr}
\toprule
Judge & $n$ & Claims seen & Title recall (raw) & Recall not in image \\
\midrule
Gemini 2.5 Flash & 50 & 0.62 & 0.20 & 0.02 \\
gpt-5 & 50 & 0.06 & 0.02 & 0.00 \\
Claude Haiku 4.5 & 50 & 0.20 & 0.10 & 0.02 \\
\bottomrule
\end{tabular}
\end{table}

\section*{S5. Pre-registration (verbatim from the project log, written before the 2,000-pair run; names redacted)}
\begin{itemize}\small
\item \textbf{Primary analysis:} prompt P1 (unchanged from src/prompts.py) on \textbf{all 2,000 pairs of pairs/clear.parquet}, for gpt-5, Claude Haiku 4.5 and Gemini 2.5 Flash, with the same settings as before. Pair credit is the same as before: 1 if the winner has the higher score, 0.5 on a tie or a failure, 0 otherwise.
\item \textbf{Comparator:} the image-blind feature baseline's out-of-fold predictions (analysis/04\_feature\_baseline\_oof\_clear.csv, computed 2026-10-08 and not refit).
\item \textbf{Equivalence test:} TOST on the paired per-pair accuracy difference (judge minus baseline) with \textbf{margin ±4 accuracy points}, α = 0.05.
\item \textbf{Robustness check:} the same analysis on pairs 501-2,000 only (the 1,500 pairs no judge had been run on before today).
\item \textbf{Frontier judge:} Claude Sonnet 5.5 (global.anthropic.claude-sonnet-5-5), P1 only, on \textbf{the first 1,000 clear pairs} (2,000 images); the budget can't cover all 2,000.
\item \textbf{Secondary analyses} (same 2,000 pairs, P1): the added-value likelihood-ratio test and the text/OCR bias probe. Bias features need OCR on the 3,000 new images.
\item \textbf{Nothing else changes:} prompts, pair sets, image preprocessing and judge settings for the three existing judges all stay as they were.
\item \textbf{Budget mechanics:} HARD\_STOP\_USD is raised to 29.20. Spend is now re-read from disk every 50 new calls, so parallel processes see each other's spend.
\item \textbf{Addition, logged before any rerun result was seen:} P1 on the 150 VisJudge-Bench calibration items for all 4 judges (about \$0.80). This lets the expert-vs-audience contrast use the same prompt and the same metric (Spearman of the P1 score with expert overall\_score vs with the audience month percentile). It addresses the review point that Spearman 0.09-0.12 against Pearson 0.41-0.47 compared different metrics and different prompts.
\item \textbf{Change to the addition above, logged before running it:} Sonnet 5.5 cost about 15\% more per image than its smoke test showed. Total spend after the rerun is about \$28.35, and P1 calibration for all 4 judges (about \$1.00) would cross the \$29.20 hard stop. P1 calibration therefore runs for gpt-5, Haiku 4.5 and Gemini 2.5 Flash only (about \$0.57). Sonnet 5.5's expert agreement is not measured.
\end{itemize}

\end{document}